%%%%%%%%%%%%%%%%%%%%%%%%%%%%%%%%%%%%%%%%%%%%%%%%%%%%%%%%%%%%%%
%%%%% MA-0781 Seminario de estudios dirigidos en matemática %%
%%%%%%%%%%%%%%%%%%%%%%%%%%%%%%%%%%%%%%%%%%%%%%%%%%%%%%%%%%%%%%
\newpage
\subsection*{MA-0781 Seminario de estudios dirigidos en matemática}
\addcontentsline{toc}{subsection}{MA-0781 Seminario de estudios dirigidos en matemática}

\hypertarget{MA0781SED}{}
\begin{refsection}
\begin{table}[H]
\centering{
\begin{tabular}{|ll|}
\hline
\textbf{Año:} V  & \textbf{Requisitos:}  \\
\textbf{Ciclo:} IX  & \textbf{Correquisitos:} Ninguno \\
\textbf{Tipo de curso:} Seminario & \textbf{Horas presenciales por semana:} 2 \\
\textbf{Créditos:} 7 & \textbf{Horas de trabajo independiente por semana:} 19 \\
\textbf{Virtualidad:} P  &  \\ \hline
\end{tabular}
}
\end{table}

\subsubsection*{Descripción del curso}

Este curso de la licenciatura ofrece un primer espacio de \textbf{estudios dirigidos} en el que la persona estudiante desarrolla, con cierta autonomía, actividades propias de la práctica investigativa en matemática, bajo la supervisión de una persona docente del Departamento de Matemáticas.

A diferencia de un curso teórico con temario fijo, aquí \textbf{los objetivos se fijan de común acuerdo} entre la persona estudiante y la persona docente al inicio del ciclo lectivo, y se consignan en un plan de trabajo. El tema puede situarse en matemática pura o aplicada, según el énfasis y los intereses de la persona estudiante, siempre que sea viable en un ciclo lectivo y pertinente a su formación.

El trabajo consiste en actividades asociadas a la investigación matemática, entre ellas: revisión documental, construcción de ejemplos y contraejemplos, demostración de teoremas, escritura de programas computacionales, ejecución de algoritmos, experimentación numérica o simbólica, y formulación o exploración de conjeturas. La persona docente \textbf{orienta} sobre los temas, las estrategias, las herramientas y otros recursos útiles para completar los objetivos; no sustituye el trabajo independiente de la persona estudiante.

El diseño se alinea con prácticas internacionales de estudio e investigación dirigida en matemática de grado avanzado (\emph{independent study}, \emph{undergraduate research}, investigación dirigida): un plan escrito al inicio, reuniones periódicas de supervisión, productos intermedios y un informe escrito con presentación oral al cierre. No se exige, como condición de aprobación, la obtención de resultados inéditos de publicación.

\subsubsection*{Objetivos}

\textbf{General:} Que la persona estudiante desarrolle, de manera autónoma, actividades asociadas a la investigación en matemática.

\textbf{Específicos:} Al finalizar el curso se espera que la persona estudiante sea capaz de:

\begin{enumerate}
\item Formular, planes de trabajo con objetivos claros, acotados y realizables.
\item Localizar, seleccionar y sintetizar literatura matemática pertinente al tema, distinguiendo resultados establecidos, técnicas disponibles y problemas abiertos o límites del enfoque adoptado.
\item Ejecutar actividades de investigación matemática acordes con los objetivos: revisión documental, construcción de ejemplos, demostración de resultados, programación, ejecución de algoritmos u otras equivalentes.
\item Elegir y utilizar, con criterio, estrategias y herramientas (analíticas, computacionales, bibliográficas u otras) recomendadas o validadas por la persona docente.
\item Documentar de forma continua el proceso de trabajo, incluyendo avances, obstáculos y ajustes al plan.
\item Comunicar de manera oral y escrita, con claridad y rigor, el contexto del tema, los métodos empleados, los resultados alcanzados y las limitaciones del trabajo.
\item Ejercer autonomía intelectual y responsabilidad académica, incluida la integridad en el uso de fuentes, datos, software y herramientas de apoyo.
\end{enumerate}

\subsubsection*{Contenidos}

Los contenidos \textbf{no se fijan de manera permanente} en el plan de estudios: dependen del tema y de los objetivos acordados para cada persona estudiante. Al inicio del ciclo lectivo deben quedar documentados por escrito:

\begin{itemize}
\item el área o tema de estudio y su justificación formativa;
\item los objetivos específicos acordados entre la persona estudiante y la persona docente;
\item las actividades previstas y un cronograma tentativo;
\item las herramientas, estrategias y fuentes iniciales de consulta;
\item los productos esperados (informe, notas, código, ejemplos, demostraciones, presentaciones u otros);
\item los criterios de seguimiento y evaluación.
\end{itemize}

A título orientativo, las actividades pueden incluir una o varias de las siguientes, según lo acordado:

\begin{enumerate}
\item \textbf{Revisión documental:} búsqueda en bases especializadas, lectura crítica de artículos o monografías, y síntesis del estado del tema.
\item \textbf{Construcción de ejemplos} y, cuando proceda, de contraejemplos que ilustren definiciones, hipótesis o la necesidad de hipótesis en un resultado.
\item \textbf{Demostración de teoremas} o de lemas auxiliares, incluyendo la adaptación de argumentos conocidos a contextos próximos.
\item \textbf{Escritura de programas computacionales} para explorar objetos matemáticos, verificar casos o apoyar conjeturas.
\item \textbf{Ejecución de algoritmos} (exactos, simbólicos o numéricos) y análisis de su comportamiento en los casos de interés.
\item \textbf{Formulación o exploración de conjeturas} a partir de evidencia teórica o computacional, con una valoración de su alcance.
\item \textbf{Comunicación de resultados:} redacción de notas o de un informe y preparación de una exposición oral.
\end{enumerate}

\subsubsection*{Metodología}

El curso adopta la modalidad de \textbf{seminario de estudios dirigidos}. El eje es el trabajo independiente de la persona estudiante (19 horas semanales), complementado con reuniones de supervisión (2 horas semanales, o su equivalente distribuido a lo largo del ciclo) en las que la persona docente orienta el proceso.

\noindent \textbf{Lineamientos metodológicos:}

\begin{enumerate}
    \item Al inicio del ciclo, la persona estudiante y la persona docente acuerdan por escrito un \textbf{plan de trabajo} con los elementos listados en Contenidos. Ese plan es el marco de las actividades y de la evaluación.
    \item Las reuniones de supervisión son periódicas y de asistencia obligatoria, salvo justificación conforme a la normativa institucional. En ellas la persona docente orienta sobre temas, estrategias, herramientas, bibliografía y ajustes de alcance.
    \item La persona estudiante mantiene un \textbf{registro de actividades} (bitácora o equivalente) y presenta avances intermedios según el calendario acordado.
    \item El alcance del plan debe corresponder a \textbf{siete créditos} en un ciclo lectivo: se privilegia un conjunto de objetivos evaluables y un producto concreto, antes que una agenda excesivamente amplia.
    \item Al finalizar, la persona estudiante entrega un \textbf{informe escrito} y realiza una \textbf{presentación oral} ante la persona docente y, si esta lo estima conveniente, ante otras personas del departamento.
    \item Se aplican criterios de integridad académica: atribución correcta de fuentes, honestidad en la descripción de resultados (incluidos los negativos o parciales) y uso responsable de software y de herramientas de apoyo, incluidas las de inteligencia artificial, cuando se autoricen en el plan.
\end{enumerate}

\textbf{Sugerencias metodológicas:}

\begin{enumerate}
    \item Acotar el tema de modo que las actividades quepan en el tiempo disponible; es preferible profundizar en pocos objetivos bien definidos.
    \item Programar al menos un hito intermedio (por ejemplo, hacia la mitad del ciclo) para revisar avances, reajustar objetivos y anticipar el informe.
    \item Recomendar el uso de bases y repositorios matemáticos (MathSciNet, zbMATH, arXiv, repositorios institucionales) y de software cuando el tema lo permita.
    \item Incorporar, si el plan lo justifica, la asistencia a coloquios o seminarios del departamento relacionados con el tema, como contexto y no como sustituto del trabajo dirigido.
    \item Redactar el informe en \LaTeX{} u otra herramienta equivalente acordada, con notación consistente y bibliografía completa.
    \item No exigir resultados originales publicables: el curso se aprueba por el cumplimiento riguroso de los objetivos acordados y por la calidad del proceso y de los productos.
\end{enumerate}

\subsubsection*{Evaluación}

La evaluación es responsabilidad de la persona docente a cargo, con base en el plan de trabajo acordado.

\subsubsection*{Bibliografía}

La bibliografía matemática del tema se define en el plan de trabajo de cada persona estudiante. Como referencias generales de escritura, búsqueda y práctica investigativa en matemática se sugieren:

\begin{enumerate}
\item Higham, N. J. (1998). \emph{Handbook of Writing for the Mathematical Sciences} (2nd ed.). SIAM.
\item Knuth, D. E., Larrabee, T., y Roberts, P. M. (1989). \emph{Mathematical Writing}. Mathematical Association of America.
\item Krantz, S. G. (2017). \emph{A Primer of Mathematical Writing} (2nd ed.). American Mathematical Society.
\item Steenrod, N. E., Halmos, P. R., Schiffer, M. M., y Dieudonné, J. A. (1973). \emph{How to Write Mathematics}. American Mathematical Society.
\item Recursos en línea de MathSciNet, zbMATH, arXiv y de la \emph{American Mathematical Society} sobre búsqueda bibliográfica y exposición matemática.
\end{enumerate}

\nocite{Kra17}

\end{refsection}
